% !TEX root = ../thesis.tex
% =============================================================================
% Chapter 7: Conclusion
% =============================================================================

\chapter{Conclusion}
\label{ch:conclusion}

This thesis has investigated whether trajectory-based behavioural features can serve as a feedback signal for LLM-driven metaheuristic design, advancing four contributions toward an answer. An extended behavioural feature vocabulary of 32 features was defined (\cref{ch:features}) and reduced via a tri-criterion consensus procedure to a tractable subset of ten predictive, non-redundant descriptors (\cref{sec:results-features}). A multi-model screening of ten LLMs identified Gemini~3 Flash as the strongest combination of generation quality, reliability, and inference cost (\cref{sec:results-screening}), with the secondary finding that medium-sized open-weight models in the 24--32B parameter range can produce viable metaheuristics. A three-way comparison of neutral, directional, and comparative feedback formats found that neutral framing, simply reporting the behavioural observation without interpretation, consistently outperformed the prescriptive variants across 29 feature-by-format conditions (\cref{sec:results-feedback}). And a head-to-head comparison on a 20-function MA-BBOB set demonstrated that all three feedback-augmented conditions, neutral-behavioural, SAGE, and combined, outperform the vanilla AOCC-only baseline (\cref{sec:results-stage4}). The decisive effect is a roughly fivefold reduction in seed-to-seed variance under combined feedback, the only Stage~4 contrast to clear $\alpha = 0.05$. The mean-AOCC lift over vanilla is directionally consistent across all three conditions but does not reach significance under Holm-corrected Welch tests at the 10-seed budget.

The answer to the central research question is therefore qualified but positive. Trajectory-based behavioural feedback does guide LLMs to design better optimisation algorithms, but the improvement manifests primarily as reduced failure-mode variance rather than as a large mean lift, and the framing of the feedback matters more than the feature content. Reporting behavioural observations to the LLM in a neutral tone helps. Telling the LLM which direction to push a feature in, even when the directional advice is empirically correct, often does not.

The first sub-question asked which LLMs can reliably generate high-quality metaheuristics for continuous black-box optimisation. The Stage~1 screening of ten models found that the closed API models lead. Gemini~3 Flash combined competitive quality ($0.862$ AOCC) with the lowest failure rate of any model ($2.0\%$) and the fastest inference, and was carried forward for that reason. The secondary finding is that medium-sized open-weight models in the 24--32B range, notably Qwen3.5~27B and Devstral Small~2, can produce viable metaheuristics, though they trail the Gemini family in both quality and reliability.

The second sub-question asked which trajectory-based behavioural features are most informative about algorithm quality. The Stage~2 tri-criterion consensus reduced the 32-feature vocabulary to ten predictive, non-redundant descriptors, with exploitation- and convergence-related features emerging as the strongest correlates of AOCC. Seven of the ten selected features were introduced to LLaMEA-generated algorithm analysis in this thesis.

The third sub-question asked whether the framing of behavioural feedback affects the performance of generated algorithms. The Stage~3 sweep across 29 feature-by-format conditions found that framing matters more than feature content. Neutral observation consistently outperformed directional and comparative prescription. Prescriptive feedback steered the median feature value in the advised direction in only $37\%$ of cases, despite that advice being empirically grounded in the Stage~1 data.

The most promising directions for future work follow directly from these caveats. Open-weight models in the 70--200B parameter range can be served on owned hardware and would lift the seed budget by an order of magnitude, bringing observed effect sizes within statistical reach. Target-model calibration of feature selection and prescriptive reference values would tighten the prescriptive signal that Stage~3 found to be blunted by behavioural-distribution collapse. Behavioural feedback should be studied at the instance or BBOB-group level before being aggregated, and should target multivariate behavioural profiles rather than individual features. AST-aware mutation guidance and retrieval-augmented prompting from a corpus of past LLaMEA-generated algorithms together offer a path toward generating algorithms that escape the gravitational pull of canonical mechanisms. None of these is a small undertaking, but each addresses a specific limitation surfaced by the experiments reported here, and together they outline a programme for moving LLM-driven algorithm design from the recombination of known mechanisms toward genuinely novel designs.
